\documentclass[11pt]{article}
\usepackage[a4paper,margin=25mm]{geometry}
\usepackage{amsmath,amssymb,amsthm,mathtools,booktabs,longtable,array}
\usepackage{mathrsfs}
\usepackage{xurl}
\usepackage{microtype}
\usepackage[colorlinks=true,linkcolor=blue,citecolor=blue,urlcolor=blue]{hyperref}
\hypersetup{hypertexnames=false,pdfauthor={}}
\setlength{\emergencystretch}{3em}
\allowdisplaybreaks
\newtheorem{theorem}{Theorem}[section]
\newtheorem{lemma}[theorem]{Lemma}
\newtheorem{proposition}[theorem]{Proposition}
\newtheorem{corollary}[theorem]{Corollary}
\theoremstyle{remark}\newtheorem{remark}[theorem]{Remark}
\newcommand{\A}{\mathcal A}
\newcommand{\B}{\mathcal B}
\newcommand{\PP}{\mathcal P_N}
\newcommand{\li}{\operatorname{li}}
\newcommand{\Th}{\Theta_N}
\newcommand{\eps}{\varepsilon}
\newcommand{\cert}[1]{\texttt{\detokenize{#1}}}
\title{A uniform seven-eighths zero-free input in an explicit Chen sieve}
\author{}
\date{7 October 2026}
\begin{document}
\maketitle
\begin{abstract}
We give an explicit reduction from a uniform zero-free half-plane for
Dirichlet $L$-functions to positive Chen representations. If every such
function is zero-free in $\Re s>7/8$, apart from the principal pole at
$1$, every even $N\ge e^{51046}$ satisfies
$D_{1,2}(N)>9\cdot10^{-7}U_NN/(\log N)^2$.
The argument uses three Chen sieve terms, exact weights at the smallest
primes, and finite Rosser weights at medium primes. We derive an effective
prime-character estimate from the zero-free input, account for small
primitive conductors even when induced to large moduli, and prove a
fixed-cutoff prime--integer rectangle estimate. Complete finite prime
enumerations and outward-rounded interval integration certify the starting
inequalities; analytic monotonicity propagates them to the whole half-line.
The external zero-free assertion is kept separate from the finite
certificates. No GRH distribution estimate is used.
\end{abstract}
\tableofcontents

\section{Result and proof plan}
Let $\Omega(n)$ count prime factors with multiplicity, and for even $N$ put
\[
D_{1,2}(N)=\#\{p<N:p\text{ prime},\ N-p\ge2,\ \Omega(N-p)\le2\}.
\]
Write $L=\log N$ and
\[
C_\infty=\prod_{p>2}\left(1-\frac1{(p-1)^2}\right),\quad
C_N=C_\infty\prod_{\substack{p>2\\p\mid N}}\frac{p-1}{p-2},\quad
U_N=2e^{-\gamma}C_N,\quad \Th=\frac{U_NN}{L^2}.
\]
The analytic input to be used is the following precise assertion:
\begin{equation}\label{eq:Z}
L(s,\chi)\ne0\quad(\Re s>7/8),
\qquad\text{except for the principal pole at }s=1.\tag{Z}
\end{equation}
It is uniform over all Dirichlet characters, moduli, and heights.
\begin{theorem}\label{thm:main}
Under \eqref{eq:Z}, every even integer $N\ge e^{51046}$ satisfies
\[
D_{1,2}(N)>9\cdot10^{-7}\Th>0.
\]
\end{theorem}
Theorem 1.1 of the external manuscript \cite{QRH} asserts exactly
\eqref{eq:Z}; adopting that theorem removes the hypothesis of
Theorem \ref{thm:main}. Its proof is not certified by our interval programs.
Our citation fixes a source revision rather than treating a downloaded
formalization entry point as a locally checked proof of its dependencies.

There are two different roles for the zero-free input. It controls a
bounded range of conductors inside the averaged ordinary AP estimate, and
it controls the low-conductor part of the switched rectangles.
High conductors are handled by published mean estimates and the
multiplicative large sieve. Summing a bound of order $x^{7/8}$ over
$\sqrt{x}$ moduli would not give the distribution estimate we need.

The combinatorial reduction has one positive sieve term and two negative
ones. Exact/finite presieving makes their distribution support inexpensive;
it replaces a large primorial with a small exact product and a bounded
medium-prime weight. The last negative term is a three-prime sequence.
Its smallest prime is separated into short intervals so that the other
two primes give a coefficient sequence bounded by one in each rectangle.
Section \ref{sec:ledger} gives the complete normalized ledger.

All parameters used in the theorem are fixed rational numbers:
\begin{equation}\label{eq:parameters}
\begin{gathered}
L_0=51046,\quad \alpha=\frac{2158047}{500000000},\quad
t=\frac12-\alpha,\quad D=N^t,\quad z=N^{1/8},\quad y=N^{1/3},\\
T=5000000,\quad w_0=5,\quad u=11/2,\quad D_0=T^u,\quad
\kappa=11/169,\quad Q_*=15,\\
\eps=1/41219,\quad \eta_+=1/3231,\quad \eta_-=1/3242,\quad
\delta=10^{-6},\quad \ell=\log(1+\delta),\quad v=10^{-8}.
\end{gathered}
\end{equation}
Here $D$ is the large-prime Rosser level, $D_0$ the defining medium level,
and $\kappa D_0$ its actual support bound. Neither is the conductor cutoff
$D_{R0}$ introduced in Section \ref{sec:rectangle}.
The certificate proves $C_\infty>0.66$ by a finite product and an integer
tail, hence throughout we use $U_N>U_*:=1.32e^{-\gamma}$.

\section{An effective consequence of the zero-free input}\label{sec:character}
We first supply the character estimate needed later. Its deliberately
generous constants avoid unspecified implied constants.
For a character $\chi$ modulo $q$, set
$\psi(x,\chi)=\sum_{n\le x}\Lambda(n)\chi(n)$ and
$\pi(x,\chi)=\sum_{p\le x}\chi(p)$.
\begin{lemma}\label{lem:character}
Assume \eqref{eq:Z}. For $q\ge1$, $x\ge2$, and every character modulo $q$,
with $a=61/64$ and $\delta_\chi=1$ for the principal character and $0$
otherwise,
\[
|\psi(x,\chi)-\delta_\chi x|\le10^6x^a\log(50qx),\qquad
|\pi(x,\chi)-\delta_\chi\li(x)|\le10^7x^a\log^2(50qx).
\]
Both bounds hold at left limits as well.
\end{lemma}
\begin{proof}
For nonprincipal $\chi$, its periodic partial sums have absolute value
at most $q$. Partial summation gives
$|L(s,\chi)|\le q|s|/\sigma$ for $\sigma=\Re s>0$.
For a principal character use the holomorphic normalization
$K(s)=(s-1)L(s,\chi)/(s+1)$, taking its removable value
$\varphi(q)/(2q)>0$ at $1$; otherwise put $K=L$.
The identity
\[
\zeta(s)=\frac{s}{s-1}-s\int_1^\infty\{r\}r^{-s-1}\,dr
\]
and the omitted Euler factors give the same growth bound for this $K$.
In detail, $|s|/|s+1|\le1$ for $\Re s>0$,
and their product has absolute value at most $2^{\omega(q)}\le q$.

On the closed disk centered at $2+i\tau$ of radius $71/64$,
$\Re s\ge57/64>7/8$, so $K$ is holomorphic and nonzero and
$|K(s)|\le10q(|\tau|+5)$. Choose $g=\log K$ using the Euler logarithm
at the center. Its absolute value there is less than $4$: the Euler
logarithm has absolute value less than $1$, and the additional principal
logarithm has absolute value at most $\log3+\pi/2$.
Put $M=\log(10q(|\tau|+5))\ge\log50$.
Borel--Carath\'eodory on the inner radius $141/128$ gives
\[
|g|\le282(M+4)+4<600M.
\]
Every point $\sigma+i\tau$ with $29/32\le\sigma\le2$ is at distance
at most $140/128$ from the center. Cauchy's formula on its radius
$1/128$ gives $|g'|<76800M$.
At $\sigma=29/32$, the extra principal logarithmic derivatives have
absolute sum at most $32/3+32/61<12$. Thus
\[
|L'/L(29/32+i\tau,\chi)|<80000\log(10q(|\tau|+5)).
\]
On horizontal segments of the strip the same bound holds for
$|\tau|\ge1$, with the pole correction at most $2$.

For $\psi_1(x,\chi)=\sum_{n\le x}(x-n)\Lambda(n)\chi(n)$,
the absolutely convergent smoothed Perron identity is
\[
\psi_1(x,\chi)=\frac1{2\pi i}\int_{(2)}
-\frac{L'}L(s,\chi)\frac{x^{s+1}}{s(s+1)}\,ds.
\]
Move the contour to $29/32$. There are no zeros in the strip;
only the principal pole contributes, with residue $x^2/2$.
The horizontal integrals tend to zero, since for fixed $x$ their
integrands are bounded by a constant times $\log(q(|\tau|+5))/\tau^2$.
The denominator on the new line is at least $.81+\tau^2$ in absolute
value. With $B=\log(50q)$,
\[
\int_0^\infty\frac{\log(10q(\tau+5))}{.81+\tau^2}\,d\tau
\le(100/81+1)(B+\log2)+1<3B.
\]
Here use $\log(1+\tau)\le\log2$ on $[0,1]$ and
$\log(1+\tau)\le\log2+\log\tau$ on $[1,\infty)$.
Since $240000/\pi<10^5$,
\[
|\psi_1(x,\chi)-\delta_\chi x^2/2|<10^5x^{61/32}\log(50q).
\]
Take $h=x^{61/64}\le x$. The forward difference quotient differs
from $\psi(x,\chi)$ by at most $(h+1)\log(x+h)$.
The two smoothed errors divided by $h$ total less than
$5\cdot10^5x^a\log(50q)$, since $1+2^{61/32}<5$.
The principal main term adds $h/2$. As $h\ge1$ and
$\log(x+h)\le\log(50qx)$, these terms imply the first bound.

Removing prime powers costs at most
$\sqrt{x}(\log x)^2/\log2\le2x^a\log^2(50qx)$.
Thus the corresponding $\theta$ error has constant $10^6+2$.
Partial summation from $2$ bounds the endpoint and integral by
$x^a\log^2(50qx)$ times
\[
\frac1{\log2}+\frac1{a\log^22}<4.
\]
The principal integration constant $2/\log2-\li(2)$ has absolute
value less than $4$, absorbed by $10^7$.
No primitivity has been assumed: the periodic sum and omitted Euler
factors use the actual modulus $q$. Left limits follow by continuity
of the error bounds.
\end{proof}

Put $r_0=3/64$ and define, for $x\ge2$,
\begin{equation}\label{eq:tail}
\tau(x)=10^7x^{-r_0}\left[
(1+r_0^{-1})\log^2(50x)+\frac{2\log(50x)}{r_0^2}+\frac2{r_0^3}\right].
\end{equation}
\begin{lemma}\label{lem:primeinputs}
For $L\ge L_0$, the following inputs hold under \eqref{eq:Z}:
\begin{gather*}
\left|\sum_{p<x}\frac1p-\log\log x-M\right|\le\tau(x),\qquad
\tau(x)<L^{-2}\quad(x\ge z),\\
N/L<|\A|<(1+2/L)N/L,\quad
\pi(x)<(1+4/L)x/\log x\quad(x\ge y),\\
\frac{U_N}{\log x}(1-v)<V_N(x)<\frac{U_N}{\log x}(1+v)
\quad(x=z,y).
\end{gather*}
Here $\A=\{N-p:p<N,\ p\nmid N\}$, $M$ is the Meissel--Mertens constant,
$\PP$ is the set of odd primes not dividing $N$, and
$V_N(x)=\prod_{p<x,p\in\PP}(1-1/(p-1))$.
If $W$ is a nonnegative monotone kernel on $[b,c]$ with $z\le b<c$,
then
\begin{equation}\label{eq:weighted}
\left|\sum_{b\le p<c}\frac{W(p)}p
-\int_b^c\frac{W(x)}{x\log x}\,dx\right|
\le 2L^{-2}\max\{W(b),W(c)\}.
\end{equation}
\end{lemma}
\begin{proof}
Partial summation of Lemma \ref{lem:character} with $q=1$ gives the
reciprocal-prime error $E(x)=e(x)/x-\int_x^\infty e(t)t^{-2}\,dt$,
where $e(x)=\pi(x)-\li(x)$. The integral converges and its bound is
exactly \eqref{eq:tail}; the constant is $M$ by taking limits.
The Stieltjes identity gives \eqref{eq:weighted}, because the sum of
endpoint contributions and total variation of $W$ is $2\max W$.

Elementary envelopes, starting at $\log x=20$, are
\[
\frac{x}{\log x}\left(1+\frac1{\log x}\right)<\li(x)
<\frac{x}{\log x}\left(1+\frac1{\log x}+\frac3{\log^2x}\right).
\]
The starting comparison is an interval evaluation of $\operatorname{Ei}(20)$;
the two derivatives minus $\li'(x)$ are respectively
$-2/\log^3x$ and $(\log x-9)/\log^4x$.
Set $e_0(b)=10^7e^{-3b/64}(b+\log50)^2$.
The three mass conclusions follow respectively from
\[
L^2e_0(L)+\frac{L^3e^{-L}}{\log2}<1,\quad
L^2e_0(L)+3/L<1,\quad
L(L/3)e_0(L/3)+27/L<1.
\]
The first pays for every prime divisor of $N$. All three conditions
are certified at $L_0$; their positive terms decrease thereafter.

The constant matches the Euler product through
\[
M=\gamma+\sum_p\left(\log(1-1/p)+1/p\right).
\]
For completeness, compare $\log\zeta(1+h)$ as $h\downarrow0$.
Partial summation of $\sum_{p\le x}1/p=\log\log x+M+o(1)$ gives
$\sum_pp^{-1-h}=\log(1/h)-\gamma+M+o(1)$, using
$\int_0^\infty e^{-u}\log u\,du=-\gamma$.
The absolutely convergent terms with prime-power exponent at least two
in $\log\zeta(1+h)$ tend to
$-\sum_p(\log(1-1/p)+1/p)$, whereas
$\log\zeta(1+h)=\log(1/h)+o(1)$. Comparing the constants proves the identity.
The logarithmic error of the Mertens product is consequently at most
$\tau(x)+1/(x-1)$, using the convergent expansion of $-\log(1-1/p)$.
The product factorization yields
\[
\frac{V_N(x)\log x}{U_N}=
e^\gamma\log x\prod_{p<x}(1-1/p)\,
\frac{\prod_{2<p<x}(1-(p-1)^{-2})}{C_\infty}
\prod_{\substack{p\mid N\\p\ge x}}\frac{p-2}{p-1}.
\]
The middle factor has logarithm between $0$ and $1/(x-2)$.
The last loses at most $2L/x$ in the present range. Both product
errors are bounded by
\[
e^{\tau(z)+4/z}-1+4L/z<v.
\]
The certificate also checks $\tau(z)<L^{-2}$. These conditions
continue to hold: their positive terms are bounded by polynomials in
$L$ times $e^{-3L/512}$ or $e^{-L/8}$, with negative derivatives
already at $L_0$. This uses no RH Mertens estimate.
\end{proof}

\section{Uniform products and finite presieving}\label{sec:presieve}
This section establishes weights whose coefficients are bounded by one
and whose support fits the later average distribution estimates.
Only the primes through $w_0$ are presieved exactly; it is not necessary
to multiply every prime through $T$ into an exact primorial.

\subsection{A product condition at every intermediate endpoint}
For $b>2$ put
\[
\mathscr D(b)=\log b\prod_{2<p<b}(1-1/(p-1)),\qquad
\mathscr D_\infty=2e^{-\gamma}C_\infty.
\]
The external unconditional inputs, for $b\ge10^4$ and in the endpoint
convention of limits, are \cite[Lemma 16 (finite upper bound) and
Lemma 17 (lower bound)]{BJS} and \cite[Lemma 3.2]{JS}:
\begin{gather*}
E(b):=\sum_{p<b}1/p-\log\log b-M\ge-2.964\cdot10^{-6}/\log b,\\
E(b)\le 2/(\sqrt b\log b)\quad(b\le10^{12}),\qquad
E(b)\le6.836\cdot10^{-6}/\log b\quad(b>10^{12}).
\end{gather*}
Let $b_A=\max\{2/(\sqrt A\log A),6.836\cdot10^{-6}/\log(10^{12})\}$
and $k_A=b_A+2.964\cdot10^{-6}/\log A+1/(A-1)+.54/A$.
For $A\ge10^4$ even, these inputs imply
\begin{equation}\label{eq:originalproduct}
\prod_{\substack{b\le p<c\\p>A}}(1-1/(p-1))^{-1}
\le e^{k_A}\frac{\log c}{\log b}\quad(1<b<c).
\end{equation}
Indeed $\log(\mathscr D(b)/\mathscr D_\infty)
=-E(b)+\sum_{p\ge b}\rho(p)$, where
$\rho(p)=-\log(1-1/(p-1))-1/p>0$ and
$\rho(p)\le1/[p(p-2)]+.54/p^2$ for $p>A$.
The two integer tails sum to at most $1/(A-1)+.54/A$.
Use the upper bound at $c$ and the lower bound at $b$.
For $b\le A<c$ take $b'\downarrow A$; for $c\le A$ the product is empty.

To sharpen this uniformly, take $S=10^8$ and $Q=10^7$, both composite,
and the finite lower bound
\[
\lambda_Q=\log2-\gamma+\sum_{2<p\le Q}\log(1-(p-1)^{-2})
+\log(1-1/Q)\le\log\mathscr D_\infty.
\]
The inequality follows from the all-integer telescoping tail
$\prod_{n>Q}(1-(n-1)^{-2})=(Q-1)/Q$.
Put $d(b)=\log\mathscr D(b)$ and bound the following quantities:
\[
A_T\ge\sup_{T\le b<c\le S}(d(b)-d(c)),\quad
M_T\ge\sup_{T\le b\le S}(d(b)-\lambda_Q),\quad
B_T=\max\{A_T,M_T+b_S,k_S\}.
\]
Finite--finite pairs are covered by $A_T$; finite--infinite pairs by
$M_T+b_S$, since $d(c)-\log\mathscr D_\infty\ge-b_S$ for $c>S$;
infinite--infinite pairs by \eqref{eq:originalproduct} at $S$.
Between primes $d$ increases, and it jumps down just after each prime.
Thus every ordered pair is covered by every left-limit peak followed by
every right-limit trough and the two composite endpoints.
For adjacent intervals their peak/trough/drawdown triples combine as
$(\max P_i,\min Q_i,\max\{A_1,A_2,P_1-Q_2\})$.
This is a complete finite algorithm, not endpoint sampling.
The certified conclusion is
\begin{equation}\label{eq:uniformproduct}
\prod_{\substack{b\le p<c\\p>T}}(1-1/(p-1))^{-1}
<(1+\eps)\frac{\log c}{\log b}\quad(1<b<c).
\end{equation}
Deletion of primes preserves it. The finite certificate enumerates all
5761454 odd primes through $S$ and proves
$e^{B_T}-1<0.000024260536<\eps$.

\subsection{Medium weights and first failure}
For any deletion of the complete primes $w_0<p\le T$, let $P_0$ be
their product and $g(p)=1/(p-1)$. For $d=p_1\cdots p_n$,
$p_1>\cdots>p_n$, set $b^\pm(d)=\mu(d)$ precisely when
\[
p_1\cdots p_m p_m^2<D_0
\]
at every odd $m\le n$ for the upper sign, and every even $m\le n$
for the lower sign; otherwise set it to zero. The empty weight is $1$.
Write $B^\pm=\sum_{d\mid P_0}b^\pm(d)g(d)$ and
$V_0=\prod_{p\mid P_0}(1-g(p))$.
\begin{lemma}\label{lem:firstfailure}
The weights bracket the roughness indicator pointwise. Their masses obey
\[
B^+=V_0+\sum_{n\text{ odd}}T_n,\qquad
B^-=V_0-\sum_{n\text{ even}}T_n,
\]
where $T_n$ sums $g(p_1\cdots p_n)\prod_{p<p_n,p\mid P_0}(1-g(p))$
over prefixes whose first checked failure is at length $n$.
Every supported divisor and every first-failing prefix has product $<D_0$.
Supported divisors in fact satisfy $d<\kappa D_0$.
\end{lemma}
\begin{proof}
Expand the full inclusion--exclusion product and partition omitted terms
by their first checked failure. For a fixed such prefix, the unrestricted
smaller-prime suffix sums to the displayed product. Its sign is positive
in $B^+-V_0$ and in $V_0-B^-$ because of the relevant parity.
The same finite identities hold at $g(p)=1_{p\mid m}$ and imply the
pointwise brackets; all these first-failure products are nonnegative.

For a supported divisor, if its last prime $r$ is checked, $dr^2<D_0$.
If it is unchecked and follows prime $s>r$, the preceding checked
prefix gives $d<D_0r/s^2$. A single unchecked prime is at most $T$.
For a first failure of length at least three, its preceding checked
prefix is two primes shorter, ending at $s>q>r$, and
$d=d_{n-2}qr<d_{n-2}s^2<D_0$.
Lengths one and two have product at most $T^2<D_0$.
This proves the product bound for every first failure.

For the sharper support use
\[
\kappa=\max\left\{p_0^{-2},\max_{w_0<r<s\le T}\frac r{s^2}\right\},
\]
with $p_0$ the smallest complete medium prime.
The checked and unchecked cases above give $d<\kappa D_0$;
also $T<\kappa T^u$ since $\kappa\ge T^{-2}$ and $u>3$.
Deletion cannot increase this complete-set maximum. For each $s$ only
its preceding complete prime need be checked. Exact integer comparisons
over all adjacent pairs give $\kappa=11/169$, attained at $(11,13)$.
\end{proof}

\subsection{A positive block bound for both masses}
This supplies a finite reproducible bound without dropping the preceding
Rosser prefix checks. Let $r_1<r_2<\cdots$ be all medium primes and let
$K$ be the largest integer with $r_1\cdots r_K<D_0$.
Take $h=1/200$ and $J=\lceil\log D_0/h\rceil$.
Label a prefix by $B=\sum\lfloor\log p/h\rfloor$ and group primes in
descending blocks of equal individual label $b$.
For each block set
\[
c=\left\lceil\frac{\log D_0-3\log p_{\min}}h\right\rceil,\qquad
f=\max\left\{0,\left\lceil\frac{\log D_0-3\log p_{\max}}h\right\rceil-K+1\right\}.
\]
For each sign maintain two nonnegative polynomials $(X_0,X_1)$,
initially $(1,0)$, for the parity of accepted prefixes.
Define $A$ to append label $b$, flip parity, and at a checked append
retain only source labels $B<c$. Truncate output labels to $B<J$.
Let $\mathcal F$ sum source coefficients $B\ge f$ of the parity checked
by one append. With $t_p=g(p)/(1-g(p))=1/(p-2)$ set
\[
D_b(a)=\prod_{p\text{ in block}}(1+t_p+t_pa),\quad
C_b(a)=\frac{D_b(a)-1}{1+a}.
\]
Both have nonnegative coefficients. Replace $X$ by
\[
\sum_{j=0}^K[a^j]D_b(a)A^jX,
\]
and add $\mathcal F\sum_{j=0}^K[a^j]C_b(a)A^jX$ to the sign's gap total.
\begin{lemma}\label{lem:block}
These two totals bound $(B^+-V_0)/V_0$ and $(V_0-B^-)/V_0$,
uniformly under every deletion of medium primes. Their certified values
satisfy
\[
(B^+-V_0)/V_0<0.000309498615<\eta_+,\qquad
(V_0-B^-)/V_0<0.000308430014<\eta_-.
\]
\end{lemma}
\begin{proof}
If a prefix has $k$ chosen primes and logarithm $v_1$, then
$hB\le v_1<h(B+k)$. Every true pass satisfies $B<c$.
A first failure has $k\le K-1$ preceding primes and
$v_1+3\log p\ge\log D_0$, which implies $B\ge f$.
All actual prefixes have product $<D_0$ by Lemma \ref{lem:firstfailure},
hence degree at most $K$ and label $<J$. The relaxed tests therefore
contain every actual pass and failure. Counting an ambiguous label both
ways only increases positive bounds.

Before processing a prime multiply the accepted-prefix mass by
$\prod_{q>p}(1-g(q))^{-1}$. One prime then gives
\[
X'=((1+t_p)I+t_pA)X,\qquad G'=G+t_p\mathcal FX.
\]
The augmented matrices commute within a block since $A,\mathcal F$ are
common. Their product gives $X'=D_b(A)X$ and
$G'=G+\mathcal FC_b(A)X$.
Equivalently $D_{\rm new}=(1+t+ta)D$ and
$C_{\rm new}=(1+t+ta)C+t$, starting from $D=1,C=0$.
If $e_m$ are elementary symmetric sums of $t_p$, then
\[
[a^j]D_b=\sum_{m\ge j}\binom mj e_m,\quad
[a^j]C_b=\sum_{m\ge j+1}\binom{m-1}j e_m.
\]
Terms $j>K$ cannot contain an actual contribution. Artificial paths
of excessive total degree may survive and harmlessly enlarge the bound.
Deletion decreases every coefficient of $D_b,C_b$; retaining the complete
blocks and their positive operators proves deletion uniformity.

The finite evaluation generates all 348510 medium primes, takes $K=22$,
$J=16968$, and processes all 1810 nonempty blocks.
Every logarithmic integer cutoff is resolved by an Arb interval with
a unique integer floor or ceiling. Polynomial coefficients use positive
Arb addition and multiplication. This evaluates exactly the finite
majorant above and proves the stated strict rational comparisons.
\end{proof}

\subsection{Composite weights and the large-prime linear sieve}
The external large-prime input is the Rosser main-mass estimate in
\cite[Theorem 6, Proposition 13 and Table 1]{BJS} with $Q=1$.
Its local density satisfies $0\le g(p)<1$, and its dimension-one product
premise holds at every intermediate endpoint by \eqref{eq:uniformproduct}.
With $s=\log D/\log x$, the masses $G_1^\pm$ satisfy
\[
G_1^-\ge V_1(x)l(s),\quad G_1^+\le V_1(x)u(s),\quad
l(s)=f(s)-108\eps e^2h(s),\quad u(s)=F(s)+106\eps e^2h(s).
\]
The lower estimate needs $D\ge x^2$ and the upper $D\ge x$.
The coefficients $\lambda_1^\pm$ are $\mu(d)$ or zero, have support
$d<D$, and their divisor sums bracket the corresponding roughness indicator.
In the ranges used here,
\[
F(s)=2e^\gamma/s\ (1\le s\le3),\quad
f(s)=2e^\gamma\log(s-1)/s\ (2\le s\le4),
\]
\[
F(s)=\frac{2e^\gamma}{s}\left(1+\int_3^s\frac{\log(x-2)}{x-1}\,dx\right)
\ (3\le s\le4),\quad
h(s)=\begin{cases}e^{-2}&1\le s\le2,\\e^{-s}&2\le s\le3,\\3e^{-s}/s&s\ge3.\end{cases}
\]
The table constants $(106,108)$ are valid since $\eps<10^{-4}$.

Let $Q_N$ contain the retained odd primes through $5$, so $Q_N\mid Q_*$,
and let $\mu_E(d)=\mu(d)1_{d\mid Q_N}$. Define
\[
\lambda^+=\mu_E*b^+*\lambda_1^+,\qquad
\lambda^-=\mu_E*(b^+*\lambda_1^--(b^+-b^-)*\lambda_1^+).
\]
For pointwise brackets $a^-\le I_0\le a^+$ and $c^-\le I_1\le c^+$,
with upper brackets nonnegative,
\[
a^-c^++a^+c^--a^+c^+\le I_0I_1\le a^+c^+.
\]
The lower expression is at most
$I_0I_1-(a^+-I_0)(c^+-I_1)$. Multiplication by exact roughness gives
the displayed composite weights.
Their prime ranges are disjoint, so factorization is unique.
Write $b^\pm(d_0)=\mu(d_0)I^\pm$ and
$\lambda_1^\pm(d_1)=\mu(d_1)J^\pm$ with indicators in $\{0,1\}$.
The lower normalized coefficient
$I^+J^--(I^+-I^-)J^+$ belongs to $\{-1,0,1\}$, by the two cases
$J^+=0,1$. Thus all full coefficients have absolute value at most one
and support $d<\kappa Q_*D_0D$.

For any nonnegative finite sequence $\mathcal C$ with
$|\mathcal C_d|=|\mathcal C|/\varphi(d)+r_{\mathcal C}(d)$,
put $R=\sum_{d<\kappa Q_*D_0D,d\mid P_N(x)}|r_{\mathcal C}(d)|$.
The main-mass calculation gives
\begin{equation}\label{eq:composite}
\begin{aligned}
S(\mathcal C,x)&\ge|\mathcal C|V_N(x)
\{(1+\eta_+)l(s)-(\eta_++\eta_-)u(s)\}-R,\\
S(\mathcal C,x)&\le|\mathcal C|V_N(x)(1+\eta_+)u(s)+R.
\end{aligned}
\end{equation}
Indeed, with $a_0=(B^+-V_0)/V_0$, $b_0=(V_0-B^-)/V_0$, the normalized
lower mass is at least $(1+a_0)l-(a_0+b_0)u$.
It decreases in $a_0,b_0$ because $u\ge l$ and $u>0$.
The positive lower-mass condition is certified for the actual $s$.
Expanding the bounded coefficients against the sequence gives $R$.

\section{Ordinary AP distribution}\label{sec:AP}
Put $x_1=N/L^5$, $w=\log x_1=L-5\log L$, and
$H=\sqrt{x_1}/\log^{10}x_1$.
The distribution cost is averaged over squarefree moduli, not estimated
one modulus at a time by Lemma \ref{lem:character}.
The external ordinary tools in \cite[Lemmas 22--24]{BJS} give
\begin{align*}
\omega(d)&<1.3841\log d/\log\log d,\\
\sum_{d\le R}\mu^2(d)&\le.65R\quad(R\ge45),\\
\sum_{d\le R}\mu^2(d)/\varphi(d)&\le1.1\log R\quad(R\ge10^9).
\end{align*}
The first bound is used for $d\ge3$.
The high-primitive-conductor input is the last displayed estimate in
the proof of \cite[Lemma 30, source label \texttt{small313}]{BJS},
obtained there from Akbary--Hambrook \cite[Theorem 1.2]{AH} and partial
summation. For $Q=\log^{10}x_1$, $x_1\le Y\le N$, it gives
\begin{equation}\label{eq:AHinput}
\log Y\sum_{\substack{Q<r\le H\\(r,N)=1}}\frac{\mu^2(r)}{\varphi(r)}
\sum_{\chi\bmod r}^{*}|\psi(Y,\chi)|\le27\mathcal E(Y),
\end{equation}
where
\[
\mathcal E(Y)=\frac{4Y\log^{9/2}Y}{\log^{10}(Y/\log^5Y)}
+\frac{4Y}{\log^{11/2}Y}
+\frac{18Y^{11/12}}{\sqrt{\log Y}}
+\frac52Y^{5/6}\log^{11/2}Y.
\]
Its required starting condition is $\log\log x_1\ge10.4$ and
$Q<H$; both are verified here. The restriction $(r,N)=1$ is retained
exactly as in the source: it holds automatically when $r\mid d$ and
$(d,N)=1$ in the application below. This high-conductor input is
unconditional and contains no zero term; the exceptional-modulus
restriction on the full BJS lemma is not used to estimate this isolated
primitive high-conductor sum. Stars mean primitive characters; in this
range none is principal.

\begin{lemma}\label{lem:AP}
Under \eqref{eq:Z}, for $L\ge L_0$ and
$r(d)=|\A_d|-|\A|/\varphi(d)$,
\[
R_A:=\sum_{\substack{d\le H\\(d,N)=1}}\mu^2(d)|r(d)|
\le c_4(N)N/L^3,\qquad c_4(N)L^2<2.21.
\]
The estimate includes all squarefree moduli, not only $d\mid P_N(z)$.
\end{lemma}
\begin{proof}
In character orthogonality write every nonprincipal character modulo
$d$ as an induction of a unique primitive character modulo $r\mid d$.
Its induction error in $\psi(Y)$ is at most
$\omega(d)\log Y\le\log^2Y/(2\log2)$.
The total such error is at most $\sqrt Y/(2\log2\log^8Y)$ since
$H\le\sqrt Y/\log^{10}Y$.
The principal deletion carries $1/\varphi(d)$, and totals less than
$.4\log^3Y$.

Exact conductor regrouping of the remaining terms is bounded by
\[
A(H)\sum_{r\le H}\frac{\mu^2(r)}{\varphi(r)}
\sum_{\chi\bmod r}^{*\prime}|\psi(Y,\chi)|,
\quad A(H)=\sum_{m\le H}\mu^2(m)/\varphi(m)\le.55\log Y.
\]
For $r\le Q$, Lemma \ref{lem:character}, the bound of $\varphi(r)$
characters, and $\log(50QY)<2\log Y$ give
$1.1\cdot10^6Y^{61/64}(\log Y)^{12}$ for the entire low-conductor sum,
including its inductions to large $d$.
At $w$, the certificate checks
$1.1\cdot10^6e^{-3w/64}w^{20}<3.2\cdot10^{-8}$.
The left side decreases for $w>1280/3$; thus this cost is at most
$3.2\cdot10^{-8}Y/\log^8Y$.
The program retains the larger allowance
$1.1\log(\log^{10}Y)3.2\cdot10^{-8}Y/\log^8Y$.
For $r>Q$, use $A(H)\le.55\log Y$ and \eqref{eq:AHinput};
the resulting cost is at most $.55\cdot27\mathcal E(Y)<27\mathcal E(Y)$.
The resulting four positive costs are exactly the zero-term-free
$\psi$ bound used in the next formula.

Convert $\psi$ to $\theta$ and then $\pi$ as in
\cite[Lemma 31, source label \texttt{small533}]{BJS}.
The prime-power bound is $1.02\sqrt Y+3Y^{1/3}$.
Brun--Titchmarsh on the initial segment $x_1$ costs $2.2x_1$.
With $E_\theta(Y;d)=\theta(Y;d,N)-\theta(Y)/\varphi(d)$, the exact
partial summation identity used is
\[
E_\pi(N;d)=E_\pi(x_1;d)+\frac{E_\theta(N;d)}L
-\frac{E_\theta(x_1;d)}w+
\int_{x_1}^{N}\frac{E_\theta(Y;d)}{Y\log^2Y}\,dY.
\]
The integral is bounded by its elementary antiderivative majorant
$N/[L^4(1-4/w)]$. Put
\begin{align*}
E_w={}&\frac{4w^{9/2}}{(w-5\log w)^{10}}+4w^{-11/2}
+18e^{-w/12}w^{-1/2}+\frac52e^{-w/6}w^{11/2},\\
p_2={}&w^2\left(1.1\cdot10\log w\frac{3.2\cdot10^{-8}}{w^8}
+27E_w+\frac{e^{-w/2}}{2\log2\,w^8}+.4w^3e^{-w}\right),\\
p_1={}&p_2+(.67+2e^{-w/6})w^{-8},\\
p={}&p_1\left(1+\frac1{L^2w^3}+\frac1{(1-4/w)L}\right)+2.2/L^2,\\
c_4(N)={}&p+.9e^{w/2-L}\frac{L^4}{w^{10}\log L}.
\end{align*}
All positive $p_2$ summands decrease as functions of $\log Y\ge w$;
therefore their supremum is the displayed endpoint.
Finally deleting the primes dividing $N$ costs at most
$1.3841L/\log L$ per modulus, and at most $.65H$ moduli occur.
This gives the last summand of $c_4$, since $.65\cdot1.3841<.9$.
This deletion argument uses only squarefreeness and $(d,N)=1$,
so the narrower printed sieve-modulus range in BJS Lemma 43 is unnecessary.
The starting inequality $c_4(N)L^2<2.21$ is certified; its half-line
extension is proved in Section \ref{sec:halfline}.
\end{proof}

\section{A fixed-cutoff rectangle estimate}\label{sec:rectangle}
The third sieve term has a prime in a short interval and an integer
coefficient formed from two larger primes. We need a uniform estimate
whose low-conductor input remains valid at its much smaller endpoints.
Put $X_3=N^{1/8}/2$, $x_3=\log X_3$, and $D_{R0}=x_3^{10}$.
For $r\le\log^{10}X_3$ and $b\ge X_3$, Lemma \ref{lem:character} gives
\begin{equation}\label{eq:v}
\sum_{\chi\bmod r,\chi\ne\chi_0}|\pi(b,\chi)|
\le v_3\frac b{\log^5b},\qquad
v_3=4\cdot10^7X_3^{-3/64}x_3^{17}.
\end{equation}
Indeed $\log(50rb)<2\log b$, there are at most $r$ characters,
and $e^{-3\xi/64}\xi^{17}$ decreases for $\xi>1088/3$.

\begin{lemma}\label{lem:rectangle}
Let $X=N/w_j$, $Z=w_j$, $Y=\min(y,(1+\delta)w_j)$, where $z\le w_j<y$.
For $|a(n)|\le1$, the sum of the absolute nonprincipal rectangle errors
over squarefree $d\le H$ with $d\mid P_N(y)$ is at most
\begin{align*}
&39v_3XY/\log^3Y+108X\log^{13}Y/\log\log Y\\
&\quad+26XY\log^2Y/D_{R0}
+88XY\log^2Y(X^{-1/2}+Y^{-1/2})+106XY/\log^9Y.
\end{align*}
The error is centered at
$\varphi(d)^{-1}\sum_{n<X}\sum_{Z\le p<Y,(np,d)=1}a(n)$.
\end{lemma}
\begin{proof}
Here $XY\ge N$, $XY\le(1+\delta)N$, $X>N^{2/3}$ and $Y>z$.
Consequently $H<\sqrt{XY}/\log^{10}Y$ and $H\le Y^4$.
Also $10^9<D_{R0}<z<H$ and $Z>X_3$.
Orthogonality and primitive conductor decomposition $d=sr$ give
\[
\sum_{sr\le H}\frac1{\varphi(sr)}\sum_{\chi\bmod r}^{*\prime}
\left|\sum_{n<X,(n,s)=1}a(n)\chi(n)\right|
\left|\sum_{Z\le p<Y,p\nmid s}\chi(p)\right|.
\]
All moduli remain squarefree; hence $(r,s)=1$ and
$\varphi(rs)=\varphi(r)\varphi(s)$.
For $r<D_{R0}$, use \eqref{eq:v} at both left limits of $Z,Y$.
Since $b/\log^5b$ increases in this range, their total is at most
$2v_3Y/\log^5Y+\varphi(r)\omega(s)$.
The integer sum is at most $X$. Two reciprocal-totient sums give
\[
2.42v_3XY\log^2H/\log^5Y
+1.21\cdot1.3841D_{R0}X\log^3H/\log\log H.
\]
In bounding $\omega(s)$ by $1.3841\log H/\log\log H$, use
monotonicity of $\log b/\log\log b$ for $b\ge16$.
For $1\le s\le15$, $\omega(s)\le3$ and the same endpoint bound
is greater than $3$, so the small values introduce no exception.
Use $H\le Y^4$, $D_{R0}\le\log^{10}Y$, and monotonicity of
$\log^3b/\log\log b$ for $b\ge5$.
The first two constants become $2.42\cdot16<39$ and
$1.21\cdot1.3841\cdot64<108$.

For high conductors split into half-open intervals
$R_j\le r<\min(2R_j,H)$, with $R_j=2^jD_{R0}$.
The number of nonempty intervals is at most $\log H/\log2$ because
$D_{R0}\ge2$, including the last interval, and $\sum_jR_j<2H$.
The multiplicative large-sieve input in the proof of
\cite[Lemma 32, source label \texttt{lemma:bf}]{BJS} is
\[
\sum_{R\le r<2R}\frac r{\varphi(r)}\sum_\chi^*
\left|\sum_{n<X}b_n\chi(n)\right|^2\le(X+12R^2)\sum_{n<X}|b_n|^2.
\]
Use it for the integer and prime sequences, whose squared norms are
at most $X$ and $Y$. Cauchy's inequality bounds one interval by
\[
XY/R+\sqrt{12X}Y+\sqrt{12Y}X+12R\sqrt{XY}.
\]
Sum over $s$ at cost $1.1\log H$, the first three terms by the number
of intervals and the fourth by their geometric sum. With $\log H\le4\log Y$
and $H\le\sqrt{XY}/\log^{10}Y$, the constants are
$1.1\cdot16/\log2<26$, $1.1\cdot16\sqrt{12}/\log2<88$, and
$1.1\cdot4\cdot24=105.6<106$. This proves every term of the estimate.
The final prime interval may be truncated; none of these bounds assumes
that it has the full relative width $\delta$.
\end{proof}

\section{Three Chen terms and their main masses}\label{sec:three}
Define $P_N(b)=\prod_{p<b,p\in\PP}p$ and
$S(\mathcal C,b)=\#\{n\in\mathcal C:(n,P_N(b))=1\}$.
For every positive $d$, let $\A_d=\{a/d:a\in\A,d\mid a\}$;
in particular this defines $\A_q$ for $q\in\PP$.
The pointwise combinatorial input is Nathanson's Theorem 10.2 as
restated in \cite[Lemma 37, source label \texttt{lem:>}]{BJS}:
\begin{equation}\label{eq:Chen}
D_{1,2}(N)\ge S(\A,z)-\frac12\sum_{z\le q<y,q\nmid N}S(\A_q,z)
-\frac12S(\B,y)-2N^{7/8}-N^{1/3}-1,
\end{equation}
where $\B=\{N-p_1p_2p_3:z\le p_1<y\le p_2\le p_3,
p_1p_2p_3<N,(p_1p_2p_3,N)=1\}$.
The extra $1$ safely excludes $N-p=1$ under either almost-prime convention.
The exponent $7/8$ here comes from a combinatorial truncation, not from
\eqref{eq:Z}. We retain it in full.

For the lower sieve use level $D$ at $z$, with $s=8t\in(3,4)$.
For the prime-factor sum use level $D/q$, giving
$s_q=8(t-\log q/L)\in[k,3]$, $k=8(t-1/3)\in(1,2)$.
For $\B$ use level $D$ at $y$, giving $3t\in(1,2)$.
The conditions $D\ge z^2$, $D/q\ge z$, $D\ge y$ follow from these ranges.
The common distribution support is legitimate because
\begin{equation}\label{eq:support}
\alpha L-\frac{25}2\log L>\log(\kappa Q_*)+u\log T,
\qquad H>\sqrt N/L^{25/2}.
\end{equation}
The first is certified at $L_0$ and increases thereafter. Thus
$\kappa Q_*D_0D<H$, and the $\A_q$ support is below $H/q$.

\subsection{The positive term}
Equation \eqref{eq:composite} and Lemma \ref{lem:primeinputs} give
\[
S(\A,z)/\Th\ge M_0-R_A/\Th,
\quad M_0=8(1-v)\{(1+\eta_+)l(s)-(\eta_++\eta_-)u(s)\}.
\]
Its bracket is strictly positive, as certified.

\subsection{The prime-factor term: a full representative calculation}
Fix $q\in[z,y)$ not dividing $N$. The sequence $\A_q$ has model
mass $|\A|/(q-1)$ and level $D/q$. For $d\mid P_N(z)$,
\[
r_q(d)=r(qd)-r(q)/\varphi(d),\qquad
|\A_q|=|\A|/(q-1)+r(q).
\]
Its main mass is bounded by
$|\A_q|V_N(z)(1+\eta_+)(F(s_q)+106\eps e^2h(s_q))$.
The kernel is nonnegative nondecreasing in $q$ and at most
$F_*=2e^\gamma/k+106\eps$.
Apply \eqref{eq:weighted}, using
$1/(q-1)\le(1-z^{-1})^{-1}/q$, and substitute $q=N^\beta$.
Since $dq/(q\log q)=d\beta/\beta$, its normalized half-main term is
\begin{equation}\label{eq:M1}
M_1=\frac{4(1+\eta_+)(1+2/L)(1+v)}{1-e^{-L/8}}
\left(I+106\eps J_h(t)+2F_*/L^2\right),
\end{equation}
where
\[
I=\frac{e^\gamma}{4t}\log\left(6\frac{3-8\alpha}{3-18\alpha}\right),
\]
\[
J_h(t)=\int_{1/8}^{t-1/4}\frac{e^{2-8(t-\beta)}}\beta\,d\beta
+\log\frac{1/3}{t-1/4}.
\]
The split is exactly $s_q=2$, where the two definitions of $h$ meet.
This derives the integral, rather than treating it as an unexplained constant.

Since $q\ge z$ and every factor of $d$ is $<z$, the map $(q,d)\mapsto qd$
is injective. Hence $\sum_q\sum_d|r(qd)|\le R_A$.
Also $\sum_q|r(q)|\le R_A$ since $y<H$, and
$\sum_{d\le H/q}\mu^2(d)/\varphi(d)\le1.1\log H<.55L$.
Here $H/q\ge10^9$ throughout. The relative sieve remainders are thus at
most $(1+.55L)R_A$. The mass-centering correction adds at most
$(1+\eta_+)V_N(z)F_*R_A$ before taking the half.
Together with the positive term, their normalized expense is bounded by
\begin{equation}\label{eq:EAP}
E_{AP}=\frac{(1.5+.275L)2.21}{U_*L^3}
+4(1+\eta_+)(1+v)F_*\frac{2.21}{L^4}.
\end{equation}
This uses the same averaged $R_A$ for every occurrence, never a GRH
pointwise bound for the prime moduli.

\subsection{The triple-prime term}
Take $w_j=z(1+\delta)^j<y$ and replace the product restriction in
each $p_1$ interval by $w_jp_2p_3<N$. Drop $(p_1,N)=1$, but retain
$(p_2p_3,N)=1$. The enlarged sequences $\overline{\B}_j$ are disjoint:
the unique prime factor below $y$ recovers $p_1$, and unique sorted
factorization of $n=p_2p_3$ recovers the remaining factors, including
$p_2=p_3$. Thus their integer coefficients are at most one.
If some differences are nonpositive, translate each sequence by a large
integer multiple of $P_N(y)$. All relevant divisibilities, masses, and
remainders are unchanged, so the nonnegative-index sieve still applies.

For $b=\log u/L$ put
\[
H_1(u)=\int_y^{\sqrt{N/u}}\frac{dv}{v\log v\log(N/(uv))}
=\frac{\log(2-3b)}{L(1-b)}.
\]
Counting the last prime by Lemma \ref{lem:primeinputs} bounds the total
mass by $(1+\delta)(1+4/L)N$ times
\[
\sum_{z\le p_1<y}\frac1{p_1}
\sum_{y\le p_2<\sqrt{(1+\delta)N/p_1}}
\frac1{p_2\log(N/(p_1p_2))}.
\]
The inner kernel is nondecreasing and at most $4/L$; its weighted
reciprocal-prime error is $8/L^3$. Extending the integral by logarithmic
length $\ell/2$ costs $6\ell/L^2$, since its logarithmic integrand is
at most $12/L^2$. The outer kernel $H_1$ is nonnegative nonincreasing
and $H_1(z)<1/L$, so its error is $2/L^3$. Finally the sum of $1/p_1$
is at most $J+2/L^2$, with $J=\log(8/3)$.
Consequently
\[
|\overline{\B}|\le\frac{(1+\delta)(1+4/L)N}{L}
\left[\bar c+2/L^2+(J+2/L^2)(6\ell/L+8/L^2)\right],
\]
where
$\bar c=\int_{1/8}^{1/3}\log(2-3b)/(b(1-b))\,db$.
Apply the upper composite sieve at $y$ and use $e^2h(3t)=1$.
The normalized half-main term is at most
\begin{equation}\label{eq:M2}
M_2=\frac{(1+\eta_+)(1+v)}2(2e^\gamma/t+318\eps)
(1+\delta)(1+4/L)
\left[\bar c+2/L^2+(J+2/L^2)(6\ell/L+8/L^2)\right].
\end{equation}

Each sequence is a rectangle with the endpoints of Lemma \ref{lem:rectangle}.
Because the integer factors are primes at least $y$, they are coprime
to every $d\mid P_N(y)$. Only deleting primes $p_1\mid d$ changes the
principal character's mass from $|\overline{\B}_j|/\varphi(d)$.
Its per-modulus difference is at most $N\omega(d)/(z\varphi(d))$.
With at most $B(L)=5L/(24\ell)+1$ boxes and
\begin{align*}
m(L)={}&39v_3+108x_3^{16}e^{-x_3}/\log x_3+26L^5/x_3^{10}\\
&+88L^5(e^{-L/3}+e^{-L/16})+106/x_3^6,
\end{align*}
the normalized half-error is bounded by
\begin{equation}\label{eq:Erect}
E_R=\frac{(1+\delta)512m(L)B(L)}{2U_*L},\qquad
E_{\rm del}=\frac{(1+\delta)1.1L^4e^{-L/8}B(L)}{2U_*}.
\end{equation}
For $E_R$, divide each term of the rectangle bound by $XY/\log^3Y$;
$\log Y\ge L/8$ gives the factor $512$, and $XY\le(1+\delta)N$.
The function $\xi^{16}e^{-\xi}/\log\xi$ decreases for $\xi\ge x_3$,
which handles its second term. For $E_{\rm del}$, the raw expense per
box is at most $1.1NL^2/z$; division by $2\Th$ supplies another $L^2$.
No factor of $L$ is discarded in this normalization.

\section{Ledger, certification, and the half-line}\label{sec:ledger}
Combining \eqref{eq:Chen}--\eqref{eq:Erect} gives
\begin{equation}\label{eq:ledger}
\frac{D_{1,2}(N)}{\Th}\ge
M_0-M_1-M_2-E_{AP}-E_R-E_{\rm del}-E_f,
\quad E_f=\frac{L^2}{U_*}(2e^{-L/8}+e^{-2L/3}+e^{-L}).
\end{equation}
This is the single ledger; no auxiliary floating-point budget enters it.
\begin{proposition}[Strict starting inequalities]\label{prop:certificate}
At $L=L_0$ with \eqref{eq:parameters}, all product, support, mass,
conductor, and sieve-range conditions in the preceding sections hold.
The right side of \eqref{eq:ledger} is strictly greater than
$9.8468\cdot10^{-7}>9\cdot10^{-7}$.
\end{proposition}
\begin{proof}
The finite algorithms and interval integrals are specified in the appendix.
For clarity, coarser decimal endpoints, all verified against their full
Arb intervals, enclose the decisive entries as follows:
\begin{center}\begin{tabular}{lll}\toprule
Entry&Lower endpoint&Upper endpoint\\\midrule
$M_0$&$7.806926624899798362$&$7.806926624899798363$\\
$M_1$&$6.500392797042414805$&$6.500392797042414806$\\
$M_2$&$1.306527039983525481$&$1.306527039983525482$\\
$E_{AP}$&$3.1474\cdot10^{-10}$&$3.1475\cdot10^{-10}$\\
$E_R$&$5.8028784809024\cdot10^{-6}$&$5.8028784809025\cdot10^{-6}$\\
$E_{\rm del}+E_f$&$0$&$10^{-2700}$\\
Net margin&$9.8468\cdot10^{-7}$&$9.8469\cdot10^{-7}$\\\bottomrule
\end{tabular}\end{center}
The support slack in \eqref{eq:support} exceeds $4.0153\cdot10^{-5}$.
The three prime-count gates in Lemma \ref{lem:primeinputs} are bounded by
$10^{-1013}$, $.000058771$, and $.000528935$, respectively.
The entire low-conductor absorption gate is below $1.328\cdot10^{-938}$.
Also $c_4(N)L^2<2.200019$, $w>50991$,
$D_{R0}>10^9$, and all the stated integer cutoffs and strict signs resolve.
The medium-prime bounds of Lemma \ref{lem:block}, ordered-product bound,
and $C_\infty>0.66$ are regenerated from complete integer prime sieves,
not obtained by reusing a GRH certificate.
The programs check the full intervals and the coarser endpoints above;
rounded decimal display alone is not used as an acceptance test.
\end{proof}

\subsection{Analytic propagation}\label{sec:halfline}
We now keep all rational parameters fixed, rather than recomputing
$\alpha$ by the search recipe. The left side of \eqref{eq:support}
increases, with derivative $\alpha-12.5/L>0$ at $L_0$ and thereafter.
Thus every distribution support continues to fit $H$.
The product condition of Section \ref{sec:presieve} already covers every
real upper endpoint, independently of $N$.

The ordinary AP envelope also propagates. Dots in the next calculation
mean derivatives with respect to $\log L$. Put
$a_1=(L-5)/w$ and $b_1=a_1(w-5)/(w-5\log w)$.
For all $L\ge L_0$,
\[
1<a_1<1.01,\quad b_1>1,\quad w>50000.
\]
For the upper bound on $a_1$, use positivity and increasingness of
$.01L-5.05\log L+5$. The other inequalities follow from $\log L>1$
and $\log w>1$.
Apart from fixed positive coefficients, the nonconstant summands in
$L^2p_1$ have forms
\[
L^2w^{-6}\log w,\quad L^2w^{13/2}(w-5\log w)^{-10},\quad
L^2w^{-7/2},\quad L^2w^{3/2}e^{-w/12},\quad
L^2w^{15/2}e^{-w/6},
\]
\[
L^2w^{-6}e^{-w/2},\quad L^2w^5e^{-w},\quad
L^2w^{-8},\quad L^2w^{-8}e^{-w/6}.
\]
Their logarithmic derivatives are respectively bounded by
$2-6+1.01/\log w$, $2+6.5\cdot1.01-10$, and $2-3.5$ for the
first three; the others have either a negative power derivative or
a negative exponential contribution $(L-5)/12$ or larger, dominating
their displayed fixed powers. They are all negative.
The correction $1+1/(L^2w^3)+1/((1-4/w)L)$ decreases.
The last summand of $L^2c_4$ is
$.9L^6e^{w/2-L}/(w^{10}\log L)$, with logarithmic derivative
$6+(L-5)/2-L-10a_1-1/\log L<0$.
Only the constant $2.2$ remains fixed. This proves $L^2c_4<2.21$
on the whole half-line, not just at its initial point.

The lower main term $M_0$ is fixed. Every varying factor in $M_1,M_2$,
namely $1+2/L$, $1+4/L$, $(1-e^{-L/8})^{-1}$, $L^{-2}$ and $\ell/L$,
decreases. Hence both negative main terms decrease.
The positive terms of $E_{AP}$ decrease individually.
For rectangles put $x_3=L/8-\log2$.
The functions $x_3^{17}e^{-3x_3/64}$,
$x_3^{16}e^{-x_3}/\log x_3$, $x_3^{-6}$, and
$L^5e^{-L/3},L^5e^{-L/16}$ all decrease at and beyond $L_0$.
The logarithmic derivative of $L^5/x_3^{10}$ is
$5-10L/(8x_3)<-5$. Also $B(L)/L=5/(24\ell)+1/L$ decreases.
Thus $E_R$ decreases. The deletion cost is a positive linear combination
of $L^5e^{-L/8}$ and $L^4e^{-L/8}$, and $E_f$ consists of decreasing
exponential-polynomial terms.
The tail and prime-mass gates in Section \ref{sec:character} decrease
by the same direct derivatives; all their exponents dominate their
fixed polynomial degrees already at $L_0$.
Every side condition persists. Discrete prime sets and the truncated
last box need not vary continuously: the common positive envelopes
already cover them. The right side of \eqref{eq:ledger} is therefore
bounded below by its certified starting value.
This proves Theorem \ref{thm:main} under its explicitly stated input.

\appendix
\section{Finite arithmetic and reproducibility}\label{sec:implementation}
The proof assigns programs only finite arithmetic tasks. Theorems about
zeros, Rosser weights, mean character sums, and the continuum propagation
are analytic dependencies; a finite program does not prove them.
The reference implementation uses Python with python-flint, Arb/ACB,
and 192-bit working precision. It runs from its own directory:
\begin{verbatim}
python qrh_chen_reproduce.py
python qrh_chen_positivity_search.py
\end{verbatim}
The first is the theorem's fresh arithmetic reproduction; the second
records a bounded parameter exploration and is not needed for the
half-line proof. All theorem parameters are declared directly as rationals.
No old GRH JSON ledger or GRH distribution evaluator is read by the
fresh reproduction.

Complete integer Eratosthenes sieves generate all primes for the medium
and ordered-product calculations. The medium block algorithm is exactly
Lemma \ref{lem:block}; multiplication of polynomials uses
\cert{arb_poly} and discards only coefficients proved irrelevant there.
The support maximum uses integer cross multiplication over adjacent
primes. The lower bound for $C_\infty$ uses the full odd prime set through
$T$, multiplied by $(T-1)/T$, the all-integer telescoping tail.
The ordered-product algorithm keeps peaks before each prime jump and
troughs after it, not sampled interior values; its analytic tail is stated
in Section \ref{sec:presieve}.

There are three interval integrations: the delayed $F(s)$ integral,
the first piece of $J_h$, and the integral $\bar c$ (the second $J_h$
piece is a logarithm). They use \cert{acb.integral}, retaining its real
enclosure. Logarithms in integrands use the integration API's analytic
flag, so branch boundaries are not silently crossed.
The path of the $F$ integral lies in $3\le x\le s<4$, with $x-2>0$;
the $J_h$ path lies in $1/8\le b\le t-1/4<1/3$ with $b>0$;
the $\bar c$ path lies in $1/8\le b\le1/3$, with $b,1-b,2-3b>0$.
Thus each integrand is holomorphic on a neighborhood of its compact
path and has no endpoint singularity. The piecewise sieve function is
split at its exact rational/parameter breakpoint before integration.
Arb's outward rounding encloses elementary evaluations and ACB's
validated integration encloses each entire integral, rather than a
floating quadrature subsequently decorated with a tolerance.

\cert{qrh_chen_reproduction.json} contains the fresh prime masses,
ordered product, two fixed rational endpoints, and source SHA-256 hashes.
\cert{qrh_chen_interface_checks.py} independently reconstructs the
deleted-prime normalization and checks the contour/rectangle constants.
\cert{qrh_chen_positivity_search.json} specifies the explored integer
range $50000\le L\le51200$, five relative widths, and the exact
$10^{-9}$ ceiling recipe for $\alpha$. It records every branch outcome;
the selected positive endpoint is 51046. The preceding recipe endpoint
is negative. This establishes neither continuous-parameter nor global
optimality. These exploration parameters are frozen, not adapted, in
Theorem \ref{thm:main}.

\section{Notation and external-input map}
\begin{longtable}{p{.22\linewidth}p{.69\linewidth}}
\toprule Symbol&Meaning\\\midrule\endhead
$L,L_0$&$\log N$ and its fixed starting value 51046.\\
$\alpha,t,D$&Support reserve, $t=1/2-\alpha$, and large Rosser level $N^t$.\\
$z,y$&Chen cutoffs $N^{1/8},N^{1/3}$.\\
$\PP,P_N(x),V_N(x)$&Odd primes not dividing $N$, their sieve product, and model density product.\\
$T,w_0,Q_*$&Medium cutoff, exact small cutoff, and complete exact odd product 15.\\
$D_0,\kappa$&Defining medium Rosser level and support contraction $11/169$.\\
$\eta_+,\eta_-$&Separate normalized medium upper/lower gap budgets.\\
$\eps,F,f,h$&Dimension-one tolerance and explicit linear-sieve functions.\\
$x_1,w,H$&$N/L^5$, its logarithm, and average AP modulus support.\\
$r(d),R_A$&Centered ordinary AP remainder and its full squarefree sum.\\
$X_3,x_3,D_{R0},v_3$&Fixed rectangle reference point, its logarithm, low-conductor cutoff, and character envelope.\\
$w_j,X,Y,Z$&Small-prime box endpoint and prime--integer rectangle bounds.\\
$U_N,U_*,\Th$&Singular-series normalization, uniform lower bound, and units of the ledger.\\
$M_0,M_1,M_2$&Positive lower main term and the two half-weighted negative main terms.\\
$E_{AP},E_R,E_{\rm del},E_f$&All ordinary, rectangle, model-deletion, and finite combinatorial charges.\\\bottomrule
\end{longtable}
The carrying external inputs are \eqref{eq:Z} from \cite[Theorem 1.1]{QRH};
large-prime Rosser masses from \cite[Theorem 6, Proposition 13, Table 1]{BJS};
the unconditional reciprocal-prime tails from \cite[Lemmas 16--17]{BJS} and
\cite[Lemma 3.2]{JS}; squarefree, totient, and divisor-count bounds from
\cite[Lemmas 22--24]{BJS}; the high-conductor estimate \eqref{eq:AHinput};
the multiplicative large-sieve inequality explicitly restated in Section
\ref{sec:rectangle}; and \eqref{eq:Chen}.
Every modification of a published sieve range is proved in the text,
not effected by dropping a printed hypothesis without a replacement.

\begin{thebibliography}{9}
\bibitem{QRH} OpenAI, \emph{The Quasi-Riemann Hypothesis: A Zero-Free
Half-Plane $\Re(s)>7/8$}, 30 September 2026, Theorem 1.1.
Source snapshot at commit
\cert{adc7f1241b42e322a6451854ab7e4b4c146bf78a}:
\url{https://github.com/openai/math/blob/adc7f1241b42e322a6451854ab7e4b4c146bf78a/preprints/The-Quasi-Riemann-Hypothesis-September-30-2026/build/paper.tex}.
\bibitem{BJS} M. Bordignon, D. R. Johnston, V. V. Starichkova,
\emph{An explicit version of Chen's theorem and the linear sieve},
Int. J. Number Theory 21 (2025), 2497--2572; version 6,
\url{https://arxiv.org/abs/2207.09452v6}.
\bibitem{AH} A. Akbary, K. Hambrook, \emph{A variant of the
Bombieri--Vinogradov theorem with explicit constants and applications},
Math. Comp. 84 (2015), 1901--1934. The displayed weighted consequence
used here is also recorded in the proof of BJS Lemma 30.
\bibitem{JS} D. R. Johnston, V. V. Starichkova,
\emph{Some explicit results on the sum of a prime and an almost prime},
\url{https://arxiv.org/abs/2208.01229}, Lemma 3.2.
\bibitem{Nathanson} M. B. Nathanson,
\emph{Additive Number Theory: The Classical Bases}, GTM 164,
Springer, 1996, Theorems 9.3 and 10.2.
\end{thebibliography}
\end{document}
